An interstellar object (ISO) enters the solar system on a hyperbolic orbit.
The distance from the Sun to the asymptote of the ISO's trajectory (impact
parameter) is $p = 1.00\,\mathrm{au}$, and the ISO has an asymptotic speed
$V_\infty = 26\,\mathrm{km\,s^{-1}}$ in the solar reference frame.

\begin{description}
\item[(a)] (5 points) Determine the perihelion distance $r_p$ of the ISO (in
  meters).
\item[(b)] (2 points) Calculate the equilibrium surface temperature $T$ of
  the object at perihelion, assuming that the ISO is spherical, fast rotating,
  and has an ideal blackbody surface composed of water ice. In space the water
  ice sublimates at $120\,\mathrm{K}$. Do we expect the water ice on ISO's
  surface to sublimate? Write YES / NO as the final answer.
\item[(c)] (3 points) Observations show that the object's apparent magnitude
  varies periodically with a light-curve peak-to-peak amplitude of
  $2.0$\,mag due to its spinning, revealing the non-spherical shape of the
  ISO. Assume that the object is an ellipsoid and the flux is proportional to
  the projected area during rotation and the ISO is rotating about its minor
  axis, which is perpendicular to the line-of-sight. Estimate the ratio of the
  major axis to minor axis of the object, $a/b$.
\end{description}

\ioaafig{2026_TH_T01-f1.pdf}
